\chapter{Finite Horizons and Factorial Shells}
\label{ch:horizonte}

% Register: D1, D2, S1. Examples first, then the definition.

\section{An afternoon in a world of 24 numbers}

Before we define, we calculate. We take three places. The first may
be 0 or 1, the second 0, 1, or 2, the third 0 to 3. The weights run
from the right: the last place counts units, the middle one counts
fours (because the last place becomes ``full'' with its four possible
values before the middle one advances), the first counts twelves.
Whoever writes down all the permitted combinations in order counts
from 0 to 23:

\begin{center}
\begin{tabular}{cccc}
$000=0$ & $001=1$ & $002=2$ & $003=3$\\
$010=4$ & $011=5$ & $012=6$ & $013=7$\\
$020=8$ & $021=9$ & $022=10$ & $023=11$\\
$100=12$ & $101=13$ & $102=14$ & $103=15$\\
$110=16$ & $111=17$ & $112=18$ & $113=19$\\
$120=20$ & $121=21$ & $122=22$ & $123=23$\\
\end{tabular}
\end{center}

No gap, no duplicate: exactly $2\cdot3\cdot4=24$ numbers. At $123$
everything stops -- every place stands at its maximum, and building
on to the left is impossible, because below base two there is
nothing. We call this world of 24 places the horizon $H_3$. Whoever
wants to reach 24 needs a new world: $H_4$ with
$2\cdot3\cdot4\cdot5=120$ places, in which all the weights are
different.

\section{The definitions}

\begin{definition}[Horizon; Register D1]
For $n\in\Nn$ let
$\Hn{n}=\{(a_1,\ldots,a_n)\mid 0\le a_i\le i\}$.
Place $i$ has base $i+1$; hence
$|\Hn{n}|=2\cdot3\cdots(n{+}1)=(n+1)!$.
\end{definition}

\begin{definition}[Value; Register D2]
$V_n(a)=\sum_{i=1}^{n}a_i\,\dfrac{(n+1)!}{(i+1)!}$, equivalently
recursively $v_0=0$ and $v_i=(i{+}1)v_{i-1}+a_i$ from left to
right, with $V_n(a)=v_n$.
\end{definition}

\begin{satz}[Register S1]
$V_n\colon\Hn{n}\to\{0,\ldots,(n+1)!-1\}$ is bijective.
\end{satz}

\begin{proof}
Induction on $n$. For $n=0$ there is nothing to show. The recursion
yields
\[
V_n(a)=(n{+}1)\,V_{n-1}(a_1,\ldots,a_{n-1})+a_n,
\qquad 0\le a_n\le n;
\]
division with remainder by $n+1$ therefore determines $a_n$
and the remaining digits uniquely (injectivity). The largest value
is $(n{+}1)(n!-1)+n=(n{+}1)!-1$, and since domain and codomain both
have $(n+1)!$ elements, surjectivity follows.
\end{proof}

The inverse map $E_n$ -- decomposing a value into its digits -- is
repeated division with remainder by $n{+}1, n, \ldots, 2$, that is,
exactly the algorithm by which one breaks seconds down into days,
hours, and minutes, only with regularly growing rather than
arbitrarily mixed bases.

\section{Factorial shells}

Every natural number $x$ has a smallest home: the
\emph{minimal horizon} $\mu(x)=\min\{n:x<(n+1)!\}$. The natural
numbers thereby split into \emph{shells} $S_n=\{x:\mu(x)=n\}$ --
and these are conceptually different from the horizons: the
horizon $\Hn{n}$ is the whole space and houses the shells
$S_0$ to $S_n$; the shell $S_n$ is its \emph{new territory}, the
values that only become possible with it.

\begin{center}
\begin{tabular}{lll}
Shell & Range & New values\\
\hline
$S_0$ & only the $0$ & $1$\\
$S_1$ & only the $1$ & $1$\\
$S_2$ & $2$ to $5$ & $4$\\
$S_3$ & $6$ to $23$ & $18$\\
$S_4$ & $24$ to $119$ & $96$\\
$S_5$ & $120$ to $719$ & $600$\\
$S_6$ & $720$ to $5039$ & $4320$\\
\end{tabular}
\end{center}

The shell boundaries are number-theoretically striking places: for
$n\ge1$ the lower boundary of the shell $S_n$ is the factorial $n!$
-- divisible by \emph{every} number from $1$ to $n$. The number
directly before it, $n!-1$, is divisible by \emph{no} number from
$2$ to $n$. The step from $719$ to $720$ thus switches, within a
single successor, from maximal coprimality to maximal divisibility
-- and at precisely this point Horimetrik forces the move into a
new world.

\begin{beispiel}
The number $7$ lives minimally in $H_3$ (since $6\le7<24$) and is
written there as $013$. In $H_4$ it is written $0012$, in $H_5$ as
$00011$, in $H_6$ as $000010$ -- and from $H_7$ onwards stably as
$0\cdots07$: as soon as the horizon is large enough, the whole
number fits into the last digit, and further leading zeros are at
last genuinely harmless. We tell this life story at leisure in an
interlude of its own later on.
\end{beispiel}

With that the foundation is laid: finite worlds of factorial size,
a bijection between digits and values, and shells with striking
boundaries. What is still missing is the true protagonist -- the
second identity of every number. It steps onto the stage in the
next chapter.
